% !TEX program = pdflatex
\documentclass[10pt,aspectratio=169]{beamer}

\usepackage[utf8]{inputenc}
\usepackage[T1]{fontenc}
\usepackage[english]{babel}
\usepackage{lmodern}
\usepackage{microtype}
\usepackage{amsmath,amssymb,bm,mathtools}
\usepackage{booktabs,array}

\usetheme[numbering=fraction,progressbar=frametitle]{metropolis}
\setbeamertemplate{frame numbering}[fraction]
\setbeamertemplate{navigation symbols}{}

\definecolor{CourseBlue}{HTML}{1A3C68}
\definecolor{CourseGreen}{HTML}{216E39}
\definecolor{CourseGold}{HTML}{B8860B}
\definecolor{CourseRed}{HTML}{9A2C2C}
\definecolor{SoftBlue}{HTML}{F2F6FA}
\definecolor{SoftGreen}{HTML}{F1F8F3}
\definecolor{SoftRed}{HTML}{FFF3F3}

\setbeamercolor{normal text}{fg=black,bg=white}
\setbeamercolor{structure}{fg=CourseBlue}
\setbeamercolor{title}{fg=CourseBlue}
\setbeamerfont{title}{size=\Large,series=\bfseries}
\setbeamercolor{frametitle}{bg=CourseBlue,fg=white}
\setbeamerfont{frametitle}{series=\bfseries,size=\large}
\setbeamertemplate{frametitle}{
  \nointerlineskip
  \begin{beamercolorbox}[wd=\paperwidth,ht=3.2ex,dp=1.4ex,leftskip=1em,rightskip=1em]{frametitle}
    \usebeamerfont{frametitle}\insertframetitle
  \end{beamercolorbox}
}
\setbeamercolor{block title}{bg=CourseBlue,fg=white}
\setbeamercolor{block body}{bg=SoftBlue,fg=black}
\setbeamercolor{block title example}{bg=CourseGreen,fg=white}
\setbeamercolor{block body example}{bg=SoftGreen,fg=black}
\setbeamercolor{block title alerted}{bg=CourseRed,fg=white}
\setbeamercolor{block body alerted}{bg=SoftRed,fg=black}
\setbeamertemplate{blocks}[rounded][shadow=false]

\newenvironment{modelbox}{\begin{beamercolorbox}[rounded=true,sep=0.7em,wd=\textwidth]{block body}
  \textbf{\color{CourseBlue}MODEL}\quad}{\end{beamercolorbox}}
\newenvironment{marketbox}{\begin{beamercolorbox}[rounded=true,sep=0.7em,wd=\textwidth]{block body example}
  \textbf{\color{CourseGreen}NO-ARBITRAGE}\quad}{\end{beamercolorbox}}
\newenvironment{simplificationbox}{\begin{beamercolorbox}[rounded=true,sep=0.7em,wd=\textwidth]{block body alerted}
  \textbf{\color{CourseRed}WATCH OUT}\quad}{\end{beamercolorbox}}
\newenvironment{pausebox}{\begin{beamercolorbox}[rounded=true,sep=0.8em,wd=\textwidth]{block body}
  \textbf{\color{CourseGold}YOUR TURN}\quad}{\end{beamercolorbox}}

\newcommand{\dd}{\,\mathrm d}
\newcommand{\E}{\mathbb E}
\newcommand{\Pp}{\mathbb P}
\newcommand{\Q}{\mathbb Q}
\newcommand{\R}{\mathbb R}
\newcommand{\cF}{\mathcal F}
\newcommand{\1}{\mathbf 1}
\newcommand{\Nlaw}{\mathcal N}
\newcommand{\esssup}{\operatorname*{ess\,sup}}
\DeclareMathOperator{\Var}{Var}
\DeclareMathOperator{\Cov}{Cov}

\AtBeginSection[]{
  \begin{frame}[plain]
    \vfill
    {\Large\bfseries\color{CourseBlue}\insertsectionhead\par}
    \medskip
    {\color{CourseBlue}\rule{0.92\textwidth}{0.8pt}}
    \vfill
  \end{frame}
}

\title{Martingale Pricing, Hedging, and No-Arbitrage}
\subtitle{From martingale prices to replication and market completeness}
\author{}
\date{Lecture 7}

\begin{document}

\begin{frame}[plain]
  \vfill
  {\LARGE\bfseries\color{CourseBlue}Martingale Pricing, Hedging, and No-Arbitrage\par}
  \medskip
  {\large From martingale prices to replication and market completeness\par}
  \medskip
  {\color{CourseGold}\rule{0.94\textwidth}{0.8pt}}
  \medskip
  {Lecture 7\par}
  \vfill
\end{frame}

\begin{frame}{One question drives the whole lecture}
\begin{block}{The question}
How do self-financing gains determine prices and hedges, and when can we
hedge every claim?
\end{block}
\begin{enumerate}
  \item Recall the pricing measure built in Lecture 6.
  \item Derive martingale valuation from self-financing gains.
  \item Replicate a dividend-paying call, then change the numeraire.
  \item Recover hedges and test whether two Brownian risks are spanned.
  \item Apply the framework to Asian claims and no-arbitrage; preview American exercise.
\end{enumerate}
\end{frame}

% ============================================================
\section{The Risk-Neutral Starting Point}
% ============================================================

\begin{frame}{Lecture 6 built the measure; Lecture 7 uses it}
\[
\boxed{
\Pp
\xrightarrow[\text{Girsanov}]{\text{change weights}}
\Q
\xrightarrow{\text{discount}}
\text{traded martingales}}
\]
\begin{block}{Today}
We start from an equivalent martingale measure \(\Q\) and ask what it implies
for prices, hedges, and market structure.
\end{block}
\begin{alertblock}{Standing caution}
A stochastic integral becomes a pricing martingale only after the required
integrability or admissibility check.
\end{alertblock}
\end{frame}


\begin{frame}{The money-market account fixes our pricing units}
For now, use the money-market account
\[
B_t^{(0)}=e^{\int_0^tr_s\dd s}.
\]
Express each traded price in these units:
\(\bar S_t^i=S_t^i/B_t^{(0)}\).
\begin{block}{Equivalent martingale measure}
\(\Q\sim\Pp\), and the discounted traded gains are
\(\Q\)-martingales when the needed integrability holds.
\end{block}
Later, a different numeraire will change both the units and the measure.
\end{frame}

\begin{frame}{Why the discounted stock is a local martingale under \(\Q\)}
For a non-dividend-paying stock under the risk-neutral measure,
\[
\dd S_t=r_tS_t\dd t+\sigma_tS_t\dd W_t^\Q,
\qquad D_t=\exp\!\left(-\int_0^t r_u\dd u\right).
\]
The product rule gives
\[
\dd(D_tS_t)=D_t\sigma_tS_t\dd W_t^\Q,
\qquad
D_tS_t=S_0+\int_0^tD_u\sigma_uS_u\dd W_u^\Q.
\]
\begin{block}{The only fact needed here}
A well-defined stochastic integral with respect to Brownian motion is a
local martingale. Extra integrability is needed to call it a true martingale.
\end{block}
\end{frame}

% ============================================================
\section{Self-Financing Leads to Martingale Prices}
% ============================================================



\begin{frame}{Discounting removes the financing account}
For now, assume the risky assets pay no dividends. Let
\[
B_t^{(0)}=\exp\!\left(\int_0^t r_u\dd u\right),
\qquad
D_t=\frac1{B_t^{(0)}}.
\]
For a self-financing wealth process
\[
X_t=\phi_t^0B_t^{(0)}+\sum_{i=1}^m\phi_t^iS_t^i,
\]
the discounted value is
\[
\bar X_t=D_tX_t
=\phi_t^0+\sum_{i=1}^m\phi_t^i\bar S_t^i.
\]
\end{frame}


\begin{frame}{Self-financing becomes a discounted stochastic integral}
The product rule gives
\[
\begin{aligned}
\dd(D_tX_t)
&=D_t\dd X_t+X_t\dd D_t\\
&=D_t\Delta_t\dd S_t
  +r_tD_t\beta_tB_t^{(0)}\dd t-r_tD_tX_t\dd t\\
&=\Delta_t\bigl(D_t\dd S_t-r_tD_tS_t\dd t\bigr).
\end{aligned}
\]
Therefore
\[
\boxed{\dd\bar X_t=\Delta_t\dd\bar S_t.}
\]
\begin{block}{Meaning}
After discounting, the bank account disappears from the gain equation.
\end{block}
\end{frame}

\begin{frame}{Replication becomes a martingale identity}
Let a self-financing strategy replicate \(H\) at \(T\), with wealth \(V_t\).
Under the required admissibility and integrability conditions,
\[
\bar V_t=\frac{V_t}{B_t^{(0)}}
\]
is a \(\Q\)-martingale and
\[
\bar V_T=\frac{H}{B_T^{(0)}}.
\]
Therefore
\[
\boxed{
V_t=B_t^{(0)}
\E^\Q\!\left[\frac{H}{B_T^{(0)}}\middle|\cF_t\right].}
\]
\end{frame}

\begin{frame}{Check bond and stock payoffs}
Assume constant \(r\) and no dividends.
\begin{columns}[T,onlytextwidth]
\column{0.48\textwidth}
\begin{block}{One unit of currency}
For \(H=1\),
\[
V_t=e^{-r(T-t)}.
\]
\end{block}
\column{0.48\textwidth}
\begin{exampleblock}{One share of stock}
For \(H=S_T\),
\[
V_t=S_t,
\]
because \(e^{-rt}S_t\) is a \(\Q\)-martingale.
\end{exampleblock}
\end{columns}
\begin{pausebox}
Verify both identities directly from the pricing formula.
\end{pausebox}
\end{frame}

\begin{frame}{Constant rates give the familiar pricing formula}
If \(r\) is constant and \(H=g(S_T)\) is replicable, then
\[
\boxed{
V_t=e^{-r(T-t)}
\E^\Q[g(S_T)\mid\cF_t].}
\]
\begin{pausebox}
Where is the physical drift \(\mu\) in this formula?
\end{pausebox}
\pause
It has been absorbed into the change from \(\Pp\) to \(\Q\).
\end{frame}

\begin{frame}{Under \(\Q\), the conditional stock law is explicit}
For constant \(r,\sigma\),
\[
\frac{\dd S_u}{S_u}=r\dd u+\sigma\dd W_u^\Q.
\]
Hence
\[
S_T=S_t\exp\!\left(
\left(r-\frac12\sigma^2\right)(T-t)
+\sigma(W_T^\Q-W_t^\Q)
\right),
\]
with \(W_T^\Q-W_t^\Q\sim\Nlaw(0,T-t)\) conditionally on \(\cF_t\).
\end{frame}

\begin{frame}{A digital payoff separates probability from price}
Let
\[
H=\1_{\{S_T>K\}}.
\]
Then
\[
V_t=e^{-r(T-t)}\Q(S_T>K\mid\cF_t).
\]
\begin{pausebox}
Is the option price itself a probability?
\end{pausebox}
\pause
No. It is a discounted risk-neutral probability.
\end{frame}

\begin{frame}{A stock-or-nothing payoff uses a weighted tail}
Let \(H=S_T\1_{\{S_T>K\}}\). Then
\[
V_t=e^{-r(T-t)}
\E^\Q[S_T\1_{\{S_T>K\}}\mid\cF_t].
\]
\begin{pausebox}
Why is this not
\(e^{-r(T-t)}S_t\Q(S_T>K\mid\cF_t)\)?
\end{pausebox}
\pause
\(S_T\) and the event are dependent. The stock-weighted tail produces the
\(d_+\) term in Black--Scholes.
\end{frame}


\begin{frame}{The forward price follows from one martingale identity}
For a forward paying \(S_T-K\) on a stock without dividends,
\[
\boxed{V_t=S_t-Ke^{-r(T-t)}.}
\]
The zero-value delivery price solves \(V_t=0\):
\[
\boxed{F_{t,T}=S_te^{r(T-t)}.}
\]
\begin{exampleblock}{Replication check}
Long one share and borrow \(Ke^{-r(T-t)}\). The terminal value is
\(S_T-K\), exactly the forward payoff.
\end{exampleblock}
\end{frame}

\begin{frame}{A call is a stock tail minus a cash tail}
For \(H=(S_T-K)^+\),
\[
(S_T-K)^+
=S_T\1_{\{S_T>K\}}-K\1_{\{S_T>K\}}.
\]
Therefore
\[
C_t=e^{-r\tau}
\E^\Q[S_T\1_{\{S_T>K\}}\mid\cF_t]
-Ke^{-r\tau}\Q(S_T>K\mid\cF_t).
\]
\begin{pausebox}
Which term becomes the \(d_+\) term and which becomes the \(d_-\) term?
\end{pausebox}
\pause
The stock-weighted tail gives \(d_+\); the cash-or-nothing tail gives \(d_-\).
\end{frame}


% ============================================================
\section{Dividend Yield and Risk-Neutral Pricing}
% ============================================================

\begin{frame}{Dividends force us to distinguish price from gain}
If one share pays continuous yield \(q_t\), its instantaneous total gain is
\[
\boxed{\dd G_t=\dd S_t+q_tS_t\dd t.}
\]
\begin{columns}[T,onlytextwidth]
\column{0.48\textwidth}
\begin{block}{Ex-dividend price}
\(S_t\) is the market value of the share after separating paid dividends.
\end{block}
\column{0.48\textwidth}
\begin{exampleblock}{Gain process}
\(G_t\) records both price changes and cash distributions.
\end{exampleblock}
\end{columns}
\begin{alertblock}{Pricing condition}
Discounted total gains, not merely the ex-dividend price, must have zero drift.
\end{alertblock}
\end{frame}


\begin{frame}{Replicate a call on a dividend-paying stock}
Under \(\Pp\), assume constant \(r,q,\sigma\), with \(\mu\) the ex-dividend drift:
\[
\dd S_t=\mu S_t\dd t+\sigma S_t\dd W_t^\Pp,
\qquad \dd B_t=rB_t\dd t.
\]
For \(V_t=v(t,S_t)\), It\^o's formula gives
\[
\dd v=(v_t+\mu S_tv_s+\tfrac12\sigma^2S_t^2v_{ss})\dd t
       +\sigma S_tv_s\dd W_t^\Pp.
\]
A self-financing portfolio with \(\Delta_t\) shares receives dividends:
\[
\boxed{\dd V_t=\Delta_t(\dd S_t+qS_t\dd t)
              +r(V_t-\Delta_tS_t)\dd t.}
\]
\end{frame}

\begin{frame}{Matching the two gains gives the dividend PDE}
Brownian exposure fixes the stock holding: \(\Delta_t=v_s(t,S_t)\).
Matching the drifts under \(\Pp\) gives
\[
v_t+\mu s v_s+\tfrac12\sigma^2s^2v_{ss}
=\mu s v_s+qs v_s+r(v-sv_s).
\]
The historical drift cancels, leaving the replication equation
\[
\boxed{v_t+\tfrac12\sigma^2s^2v_{ss}+(r-q)sv_s-rv=0,
\qquad v(T,s)=(s-K)^+.}
\]
\begin{marketbox}
The call costs \(v(t,S_t)\): the stock and bank-account strategy replicate
its terminal payoff.
\end{marketbox}
\end{frame}

\begin{frame}{A prepaid share explains every dividend adjustment}
A claim delivering one share at \(T\) is worth
\[
\boxed{S_te^{-q(T-t)}}.
\]
\begin{pausebox}
How many shares should we buy at time \(t\) if all dividends are reinvested?
\end{pausebox}
\pause
Buy \(e^{-q(T-t)}\) shares. Reinvestment grows the position at rate \(q\),
so it reaches one share at \(T\).
\[
F_{t,T}=S_te^{(r-q)(T-t)}.
\]
\end{frame}

\begin{frame}{Recover the risk-neutral measure with a dividend yield}
Under \(\Pp\), let \(\mu\) be the \emph{ex-dividend price drift}:
\[
\frac{\dd S_t}{S_t}=\mu\dd t+\sigma\dd W_t^\Pp,
\qquad \text{dividend flow}=qS_t\dd t.
\]
Choose \(\lambda=(\mu+q-r)/\sigma\) and set
\[
Z_T=e^{-\lambda W_T^\Pp-\lambda^2T/2},
\qquad W_t^\Q=W_t^\Pp+\lambda t.
\]
\begin{pausebox}
Substitute \(\dd W_t^\Pp=\dd W_t^\Q-\lambda\dd t\). What is the new
ex-dividend drift?
\end{pausebox}
\pause
\[\boxed{\dd S_t/S_t=(r-q)\dd t+\sigma\dd W_t^\Q.}\]
\end{frame}

\begin{frame}{Solve the replication PDE with two Gaussian tails}
The PDE's Feynman--Kac solution, with \(\tau=T-t\), is
\[
v(t,s)=e^{-r\tau}\E^\Q[(S_T-K)^+\mid S_t=s].
\]
Here
\[
S_T=s\exp\!\left((r-q-\tfrac12\sigma^2)\tau+\sigma\sqrt\tau Z\right),
\quad Z\sim\Nlaw(0,1).
\]
Define
\[
d_2^q=\frac{\log(s/K)+(r-q-\sigma^2/2)\tau}{\sigma\sqrt\tau},
\qquad d_1^q=d_2^q+\sigma\sqrt\tau.
\]
The cash and stock tails are
\[
\Q(S_T>K\mid S_t=s)=\Phi(d_2^q),
\qquad
e^{-r\tau}\E^\Q[S_T\1_{\{S_T>K\}}\mid S_t=s]
=se^{-q\tau}\Phi(d_1^q).
\]
\end{frame}

\begin{frame}{Replication fixes the dividend-adjusted call price and hedge}
Using \((S_T-K)^+=S_T\1_{\{S_T>K\}}-K\1_{\{S_T>K\}}\),
\[
\boxed{C_t=S_te^{-q\tau}\Phi(d_1^q)-Ke^{-r\tau}\Phi(d_2^q).}
\]
Differentiating the solution gives the replicating stock holding
\[
\boxed{\Delta_t=v_s(t,S_t)=e^{-q\tau}\Phi(d_1^q),}
\qquad
C_t-\Delta_tS_t=-Ke^{-r\tau}\Phi(d_2^q).
\]
\begin{marketbox}
This value solves the replication PDE and has the call payoff at maturity;
the stock and bank-account holdings therefore fix its no-arbitrage price.
\end{marketbox}
\end{frame}

% ============================================================
\section{Changing the Numeraire}
% ============================================================

\begin{frame}{A general numeraire changes both units and measure}
Let \(N_t>0\) be a traded numeraire. Under its associated measure \(\Q^N\),
traded prices expressed in units of \(N\) are martingales:
\[
\frac{X_t}{N_t}
=\E^{\Q^N}\!\left[\frac{X_T}{N_T}\middle|\cF_t\right].
\]
Thus a replicable payoff \(H\) has value
\[
\boxed{
V_t=N_t\E^{\Q^N}\!\left[\frac{H}{N_T}\middle|\cF_t\right].}
\]
\begin{block}{Why choose another numeraire?}
A suitable unit can turn a complicated discounted payoff into a simpler
conditional expectation.
\end{block}
\end{frame}

\begin{frame}{The best numeraire can make a payoff trivial}
Suppose \(H=N_TY_T\). Then
\[
\frac{V_t}{N_t}
=\E^{\Q^N}[Y_T\mid\cF_t].
\]
\begin{pausebox}
What happens if \(H=N_T\)? What is the price without any calculation?
\end{pausebox}
\pause
\[
Y_T=1
\quad\Longrightarrow\quad
\boxed{V_t=N_t.}
\]
Pricing becomes transparent when the payoff is naturally expressed in the
chosen unit.
\end{frame}

% ============================================================
\section{From Prices to Hedges: Martingale Representation}
% ============================================================

\begin{frame}{Start from a terminal variable and its conditional expectation}
Let \(\xi\in L^2(\Q)\) be a terminal random variable and define
\[
\boxed{M_t:=\E^\Q[\xi\mid\cF_t].}
\]
Then \(M\) is a square-integrable martingale, with \(M_T=\xi\).
\begin{block}{What is missing?}
The conditional expectation identifies the martingale. How do we express
its random increments as a stochastic integral?
\end{block}
\begin{exampleblock}{The route}
Terminal variable \(\xi\) \(\longrightarrow\) martingale \(M\)
\(\longrightarrow\) Brownian coefficient \(\psi\)
\(\longrightarrow\) hedge holdings.
\end{exampleblock}
\end{frame}

\begin{frame}{The MRT decomposes every Brownian martingale into shocks}
Assume \((\cF_t)\) is the augmented filtration of a \(d\)-dimensional
\(\Q\)-Brownian motion \(W^\Q\), and \(\xi\in L^2(\Q)\).
\begin{block}{Martingale Representation Theorem}
There exists a predictable process \(\psi_t\in\R^d\) such that
\[
\boxed{
M_t=M_0+\int_0^t\psi_u^\top\dd W_u^\Q,}
\]
with \(\E^\Q[\int_0^T\|\psi_u\|^2\dd u]<\infty\).
\end{block}
\begin{exampleblock}{Interpretation}
\(\psi_t\) is the martingale's instantaneous exposure to each Brownian risk.
\end{exampleblock}
\end{frame}

\begin{frame}{Example: the representation of a linear terminal shock}
Take \(d=1\) and \(\xi=W_T^\Q\). Write
\[
W_T^\Q=W_t^\Q+(W_T^\Q-W_t^\Q).
\]
The increment is independent of \(\cF_t\) and has mean zero, so
\[
M_t=\E^\Q[W_T^\Q\mid\cF_t]=W_t^\Q.
\]
\[
\boxed{M_t=M_0+\int_0^t1\dd W_u^\Q,}
\qquad
\boxed{\psi_t=1.}
\]
\begin{block}{Reading}
The claim carries one unit of Brownian exposure at every time.
\end{block}
\end{frame}

\begin{frame}{Example: a quadratic claim has state-dependent exposure}
Still with \(d=1\), take \(\xi=(W_T^\Q)^2\). Expanding the square gives
\[
\begin{aligned}
M_t
&=(W_t^\Q)^2
 +2W_t^\Q\underbrace{\E^\Q[W_T^\Q-W_t^\Q\mid\cF_t]}_{0}\\
&\quad+\underbrace{\E^\Q[(W_T^\Q-W_t^\Q)^2\mid\cF_t]}_{T-t}
 =(W_t^\Q)^2+(T-t).
\end{aligned}
\]
It\^o's formula shows the cancellation:
\[
\dd M_t=2W_t^\Q\dd W_t^\Q+\dd t-\dd t
=2W_t^\Q\dd W_t^\Q.
\]
\[
\boxed{M_t=T+\int_0^t2W_u^\Q\dd W_u^\Q,}
\qquad \boxed{\psi_t=2W_t^\Q.}
\]
\end{frame}


\begin{frame}{The discounted stock has its own Brownian exposure}
For the hedging application, work directly with discounted values.
With one stock, one Brownian factor, and no dividends, write
\[
\boxed{\dd\bar S_t=\sigma_t\bar S_t\dd W_t^\Q.}
\]
The terminal variable is now \(\xi=\bar V_T\), so \(M_t=\bar V_t\).
If a self-financing hedge holds \(\Delta_t\) shares, then
\[
\dd\bar V_t=\Delta_t\dd\bar S_t
=\Delta_t\sigma_t\bar S_t\dd W_t^\Q.
\]
\end{frame}

\begin{frame}{Matching Brownian coefficients reveals the hedge}
The MRT gives
\[
\dd\bar V_t=\psi_t\dd W_t^\Q,
\]
while self-financing gives
\[
\dd\bar V_t=\Delta_t\sigma_t\bar S_t\dd W_t^\Q.
\]
If \(\sigma_t\bar S_t\neq0\), matching the coefficients gives
\[
\boxed{
\Delta_t=\frac{\psi_t}{\sigma_t\bar S_t}.}
\]
\begin{block}{Central message}
The martingale approach does not lose the hedge. It stores it inside
\(\psi_t\).
\end{block}
\end{frame}

\begin{frame}{In a smooth Markov model the MRT hedge is the usual delta}
Suppose \(V_t=v(t,S_t)\). The Brownian part of the discounted value is
\[
\dd\bar V_t=v_s(t,S_t)\sigma_t\bar S_t\dd W_t^\Q.
\]
Thus the MRT coefficient is
\[
\psi_t=\sigma_t\bar S_t v_s(t,S_t).
\]
Substitution into the MRT hedge formula gives
\[
\boxed{\Delta_t=\frac{\sigma_t\bar S_t v_s(t,S_t)}
{\sigma_t\bar S_t}=v_s(t,S_t).}
\]
\begin{marketbox}
The PDE delta and the martingale-representation hedge are the same object.
\end{marketbox}
\end{frame}

% ============================================================
\section{Two Brownian Risks and Market Completeness}
% ============================================================

\begin{frame}{Begin again under the historical measure}
Think of \(W^1\) as an AI-sector shock and \(W^2\) as a broad market shock.
Under \(\Pp\), model an AI stock \(S\) and a sector state \(Y\) by
\[
\frac{\dd S_t}{S_t}
=\mu(Y_t)\dd t+\sigma_1(Y_t)\dd W_t^1+\sigma_2\dd W_t^2,
\qquad
\dd Y_t=a(Y_t)\dd t+b(Y_t)\dd W_t^1.
\]
Assume \(\sigma_2\neq0\). The sector state affects stock volatility,
but \(Y\) is not traded.
\begin{pausebox}
We assumed \(W^1\) and \(W^2\) independent. Is that realistic or useful?
What might change if sector and market shocks are correlated?
\end{pausebox}
\pause
Independence simplifies the calculation; observed sector and market shocks
may be correlated.
\end{frame}

\begin{frame}{One traded stock does not hedge every call exposure}
For a call \(H=(S_T-K)^+\), write \(V_t=v(t,S_t,Y_t)\).
Its two Brownian coefficients from It\^o's formula are
\[
\underbrace{\bigl(S_t\sigma_1(Y_t)v_s+b(Y_t)v_y,\;
                   S_t\sigma_2v_s\bigr)}_{\text{call exposure}}.
\]
Holding \(\Delta_t\) shares gives only
\[
\underbrace{\bigl(\Delta_tS_t\sigma_1(Y_t),\;
                   \Delta_tS_t\sigma_2\bigr)}_{\text{stock exposure}}.
\]
\begin{alertblock}{Where replication fails}
The market component forces \(\Delta_t=v_s\), leaving the sector exposure
\(b(Y_t)v_y\) uncovered when it is nonzero.
\end{alertblock}
\end{frame}

\begin{frame}{Why the extra state is essential to this example}
If \(S\) had constant volatility and the payoff depended only on \(S_T\),
then \(v(t,S_t)\) would have \(v_y=0\): ordinary delta hedging would still
replicate the call, even with two Brownian terms in \(S\).
\begin{block}{The actual obstruction}
Here sector conditions change the call value beyond today's stock price.
That produces a second, independent hedge exposure through \(v_y\).
\end{block}
\begin{pausebox}
What additional traded instrument could supply exposure to the AI-sector
shock?
\end{pausebox}
\end{frame}

\begin{frame}{A second traded asset can span the missing shock}
Suppose another traded asset \(R\) has its own drift and both exposures:
\[
\frac{\dd R_t}{R_t}
=\nu(Y_t)\dd t+\eta_1(Y_t)\dd W_t^1+\eta_2(Y_t)\dd W_t^2.
\]
The stock and \(R\) can match the call's stochastic terms when
\[
\begin{pmatrix}
S\sigma_1 & R\eta_1\\ S\sigma_2 & R\eta_2
\end{pmatrix}
\begin{pmatrix}\Delta^S\\\Delta^R\end{pmatrix}
=\begin{pmatrix}S\sigma_1v_s+bv_y\\S\sigma_2v_s\end{pmatrix}.
\]
\begin{block}{Rank condition}
The hedge is unique when \(\sigma_1\eta_2-\sigma_2\eta_1\neq0\): the two
traded assets cover two independent shock directions.
\end{block}
\end{frame}

\begin{frame}{This spanning property is market completeness}
A market is \emph{complete} when every suitable terminal payoff in the
modeled filtration can be replicated by self-financing trading.
\[
\boxed{\text{two Brownian risks}
\quad+\quad\text{two independent traded exposures}
\quad\Longrightarrow\quad\text{complete market}}
\]
The rank criterion applies under the usual admissibility and integrability
conditions.
\begin{columns}[T,onlytextwidth]
\column{0.48\textwidth}
\begin{exampleblock}{Full rank}
Both risk directions are spanned; claims have unique replication prices.
\end{exampleblock}
\column{0.48\textwidth}
\begin{alertblock}{Rank one}
Some risks remain unhedged; no-arbitrage may leave a range of prices.
\end{alertblock}
\end{columns}
\end{frame}

\begin{frame}{The martingale route reaches the same hedge}
With full rank, solve for the market prices of risk \(\lambda_1,\lambda_2\):
\[
\begin{pmatrix}\sigma_1&\sigma_2\\\eta_1&\eta_2\end{pmatrix}
\begin{pmatrix}\lambda_1\\\lambda_2\end{pmatrix}
=\begin{pmatrix}\mu-r\\\nu-r\end{pmatrix}.
\]
Assuming the change-of-measure density is valid,
\(\dd W_t^{i,\Q}=\dd W_t^{i,\Pp}+\lambda_i(t)\dd t\), and both discounted
asset prices are \(\Q\)-local martingales.
\[
\boxed{V_t=e^{-r(T-t)}\E^\Q[(S_T-K)^+\mid\cF_t].}
\]
Martingale representation gives two shock coefficients; matching them to
the two traded assets recovers \((\Delta^S,\Delta^R)\).
\end{frame}

\begin{frame}{With several assets, hedging becomes a linear system}
Let \(m\) discounted traded assets be driven by \(d\) Brownian factors:
\[
\dd\bar S_t=\Gamma_t\dd W_t^\Q,
\qquad
\Gamma_t\in\R^{m\times d}.
\]
If
\[
\dd\bar V_t=\psi_t^\top\dd W_t^\Q
\]
and the hedge uses holdings \(\phi_t\in\R^m\), then
\[
\dd\bar V_t=\phi_t^\top\Gamma_t\dd W_t^\Q.
\]
Hence
\[
\boxed{\phi_t^\top\Gamma_t=\psi_t^\top.}
\]
\end{frame}


\begin{frame}{Unspanned risk creates several martingale measures}
Suppose the stock is driven only by \(W^1\):
\[
\frac{\dd S_t}{S_t}=\mu\dd t+\sigma\dd W_t^1.
\]
Girsanov must choose
\[
\lambda=\frac{\mu-r}{\sigma}
\]
for \(W^1\), but it may shift the independent \(W^2\) by any \(\eta\).
\[
Z_T^\eta
=\exp\!\left(
-\lambda W_T^1-\eta W_T^2
-\tfrac12(\lambda^2+\eta^2)T\right).
\]
Every \(\eta\) prices the traded stock consistently.
\end{frame}

\begin{frame}{A nontraded claim exposes the price ambiguity}
Consider
\[
H=(W_T^2)^2.
\]
Under \(\Q^\eta\),
\[
W_T^2\sim\Nlaw(-\eta T,T),
\]
so
\[
\boxed{
\E^{\Q^\eta}[H]=T+\eta^2T^2.}
\]
\begin{pausebox}
Why do different expectation values not create an arbitrage?
\end{pausebox}
\pause
Because \(H\) is not replicable by assets driven only by \(W^1\).
\end{frame}

\begin{frame}{Rank connects Brownian risks, hedges, and completeness}
For
\[
\dd\bar S_t=\Gamma_t\dd W_t^\Q,
\]
completeness requires every \(\psi_t\in\R^d\) to be attainable as
\(\phi_t^\top\Gamma_t\).
\begin{block}{Linear-algebra criterion}
The traded volatility matrix must span all Brownian directions:
\[
\boxed{\operatorname{rank}(\Gamma_t)=d.}
\]
\end{block}
This requires at least as many independent traded exposures as Brownian risks.
\end{frame}

\begin{frame}{A redundant asset can create many hedges without incompleteness}
Let
\[
\Gamma=
\begin{pmatrix}
1\\2
\end{pmatrix},
\qquad d=1,\quad m=2.
\]
For claim exposure \(\psi=3\), hedge holdings satisfy
\[
\phi_1+2\phi_2=3.
\]
\begin{pausebox}
Give two different hedges. Does nonuniqueness of holdings imply nonuniqueness
of the claim price?
\end{pausebox}
\pause
Examples are \((3,0)\) and \((1,1)\). The risk is fully spanned; the second
asset is redundant.
\end{frame}

\begin{frame}{Completeness concerns payoffs, not portfolio labels}
\begin{columns}[T,onlytextwidth]
\column{0.48\textwidth}
\begin{block}{Complete market}
Every admissible Brownian exposure is spanned. Replicable claims have a
unique no-arbitrage price.
\end{block}
\column{0.48\textwidth}
\begin{alertblock}{Incomplete market}
Some exposures are unspanned. Several EMMs may agree on traded assets but
disagree on nontraded claims.
\end{alertblock}
\end{columns}
\begin{exampleblock}{Do not confuse}
Many portfolio representations may reflect redundant assets; many pricing
measures reflect missing risk-spanning assets.
\end{exampleblock}
\end{frame}

% ============================================================
\section{Asian Options and State Augmentation}
% ============================================================

\begin{frame}{Markov structure simplifies notation, not the pricing principle}
\begin{columns}[T,onlytextwidth]
\column{0.48\textwidth}
\begin{block}{Markov payoff}
\[
H=g(S_T),\qquad V_t=v(t,S_t).
\]
The expectation may lead to a PDE or a closed form.
\end{block}
\column{0.48\textwidth}
\begin{exampleblock}{Path-dependent payoff}
\[
H=G((S_u)_{0\leq u\leq T}).
\]
The same conditional expectation remains valid on the full filtration.
\end{exampleblock}
\end{columns}
\begin{block}{Key distinction}
The pricing identity is general; the low-dimensional state representation is
the special feature.
\end{block}
\end{frame}

\begin{frame}{A complete market can still have path-dependent payoffs}
Return to the complete one-factor Black--Scholes market. A fixed-strike
arithmetic Asian call pays
\[
H=\left(\frac1T\int_0^T S_u\dd u-K\right)^+.
\]
\begin{block}{Why \(S_T\) is insufficient}
Two paths can end at the same stock price while having very different running
averages and therefore different payoffs.
\end{block}
\end{frame}


\begin{frame}{A running integral restores a Markov state}
Define
\[
\boxed{A_t=\int_0^tS_u\dd u.}
\]
Then
\[
\dd A_t=S_t\dd t,
\qquad
H=\left(\frac{A_T}{T}-K\right)^+.
\]
If \(S\) is Markov, the pair \((S_t,A_t)\) is Markov.
\begin{marketbox}
Path dependence often becomes ordinary state dependence after adding the
right accumulated variable.
\end{marketbox}
\end{frame}

\begin{frame}{The pricing identity survives state enlargement unchanged}
Under constant \(r\),
\[
\boxed{
V_t=e^{-r(T-t)}
\E^\Q\!\left[
\left(\frac{A_T}{T}-K\right)^+
\middle|\cF_t\right].}
\]
In a Markov specification,
\[
V_t=v(t,S_t,A_t).
\]
\begin{block}{What changed?}
Not the expectation formula. Only the state required to evaluate it.
\end{block}
\end{frame}


\begin{frame}{The running average adds a transport term}
Under \(\Q\), \(\dd S_t=rS_t\dd t+\sigma S_t\dd W_t^\Q\) and
\(\dd A_t=S_t\dd t\). It\^o's formula for \(v(t,S_t,A_t)\) gives
\[
\dd v=
\left(v_t+rs\,v_s+s\,v_a+\tfrac12\sigma^2s^2v_{ss}\right)\dd t
+\sigma s\,v_s\dd W_t^\Q.
\]
Therefore
\[
\boxed{
v_t+rs\,v_s+s\,v_a
+\tfrac12\sigma^2s^2v_{ss}-rv=0,}
\]
with
\[
v(T,s,a)=\left(\frac aT-K\right)^+.
\]
\end{frame}

\begin{frame}{Arithmetic and geometric averages behave differently}
\begin{columns}[T,onlytextwidth]
\column{0.48\textwidth}
\begin{block}{Arithmetic average}
\[
\frac1T\int_0^T S_u\dd u.
\]
Common in contracts, but generally no simple Black--Scholes closed form.
\end{block}
\column{0.48\textwidth}
\begin{exampleblock}{Geometric average}
\[
\exp\!\left(\frac1T\int_0^T\log S_u\dd u\right).
\]
Its logarithm is Gaussian in Black--Scholes, making it tractable.
\end{exampleblock}
\end{columns}
\end{frame}


\begin{frame}{Monte Carlo handles path dependence without a closed form}
For each simulated path under \(\Q\):
\begin{enumerate}
  \item generate \(S_{t_1},\ldots,S_{t_n}\);
  \item compute the path average;
  \item evaluate the payoff;
  \item discount and average across paths.
\end{enumerate}
\[
\widehat V_0
=e^{-rT}\frac1M\sum_{m=1}^M
\left(\frac1n\sum_{k=1}^nS_{t_k}^{(m)}-K\right)^+.
\]
\begin{block}{Strength}
Simulation uses the martingale pricing identity directly, even when the PDE
state is high-dimensional.
\end{block}
\end{frame}

%% ============================================================
%\section{No-Arbitrage Principles and the FTAP}
%% ============================================================
%
%\begin{frame}{Why does a pricing measure exist, and when is it unique?}
%We have priced attainable claims under an equivalent martingale measure.
%The fundamental theorems connect that measure to market structure:
%\begin{itemize}
%  \item no arbitrage supports the existence of a pricing measure;
%  \item spanning all modeled risks makes it unique.
%\end{itemize}
%We begin with the elementary dominance and law-of-one-price arguments.
%\end{frame}
%
%\begin{frame}{Dominance is the fastest no-arbitrage comparison}
%Let two claims have payoffs \(A\) and \(B\) at the same maturity.
%\[
%A\geq B\quad\text{in every state}.
%\]
%No-arbitrage requires
%\[
%\boxed{\operatorname{Price}(A)\geq\operatorname{Price}(B).}
%\]
%\begin{marketbox}
%A larger payoff cannot have a smaller price.
%\end{marketbox}
%\end{frame}
%
%\begin{frame}{Can you construct the arbitrage when dominance is violated?}
%Suppose \(A\geq B\) but \(a<b\), where \(a,b\) are current prices.
%\begin{pausebox}
%Which claim do you buy and which do you sell? Track cash at time \(0\) and
%payoff at \(T\).
%\end{pausebox}
%\pause
%Buy \(A\) and sell \(B\):
%\[
%\text{initial receipt}=b-a>0,
%\qquad
%\text{terminal payoff}=A-B\geq0.
%\]
%\end{frame}
%
%\begin{frame}{The law of one price is dominance in both directions}
%If two strategies satisfy
%\[
%X_T=Y_T\quad\text{in every state},
%\]
%then
%\[
%\boxed{X_0=Y_0.}
%\]
%Otherwise buy the cheaper strategy and sell the more expensive one.
%\begin{block}{Derivative pricing consequence}
%If a payoff is exactly replicated, its price must equal the replication cost.
%\end{block}
%\end{frame}
%
%
%\begin{frame}{An EMM turns an economic restriction into probability}
%A measure \(\Q\) is an equivalent martingale measure when
%\begin{enumerate}
%  \item \(\Q\sim\Pp\): both measures have the same null events;
%  \item discounted traded asset prices are \(\Q\)-martingales
%        (or local martingales in a broader formulation).
%\end{enumerate}
%\begin{block}{Equivalence}
%Possible states remain possible. Only their probability weights change.
%\end{block}
%\end{frame}
%
%\begin{frame}{FTAP I links no-arbitrage to existence of an EMM}
%\begin{block}{First Fundamental Theorem, informal}
%A market is arbitrage-free precisely when an equivalent martingale measure
%exists, subject to the theorem's technical framework.
%\end{block}
%\[
%\boxed{
%\text{no arbitrage}
%\quad\Longleftrightarrow\quad
%\text{existence of an EMM}.}
%\]
%\begin{itemize}
%  \item Left side: an economic statement about trading strategies.
%  \item Right side: a probabilistic statement about discounted prices.
%\end{itemize}
%\end{frame}
%
%\begin{frame}{FTAP II links completeness to uniqueness of the EMM}
%\begin{block}{Second Fundamental Theorem, informal}
%In an arbitrage-free market,
%\[
%\boxed{
%\text{market complete}
%\quad\Longleftrightarrow\quad
%\text{EMM unique}.}
%\]
%\end{block}
%\begin{itemize}
%  \item Completeness means every admissible claim can be replicated.
%  \item Uniqueness means all no-arbitrage pricing rules agree.
%\end{itemize}
%\end{frame}
%
%\begin{frame}{Existence and uniqueness answer different questions}
%\begin{center}
%\begin{tabular}{p{0.28\textwidth}p{0.30\textwidth}p{0.30\textwidth}}
%\toprule
%Question & Market statement & Measure statement\\
%\midrule
%Is arbitrage absent? & no free lunch & an EMM exists\\
%Is every risk spanned? & completeness & the EMM is unique\\
%Is this claim unique? & claim replicable & all EMM expectations agree\\
%\bottomrule
%\end{tabular}
%\end{center}
%\begin{simplificationbox}
%A unique price for one particular claim does not by itself prove the whole
%market is complete.
%\end{simplificationbox}
%\end{frame}
%
%\begin{frame}{Superhedging produces a rigorous upper price bound}
%Suppose a self-financing strategy satisfies
%\[
%X_T\geq H
%\quad\text{in every state}.
%\]
%Then no-arbitrage requires
%\[
%\boxed{V_0\leq X_0.}
%\]
%Similarly, if \(Y_T\leq H\), then
%\[
%\boxed{Y_0\leq V_0.}
%\]
%\begin{block}{Exact replication}
%When \(X_T=H\), upper and lower bounds meet at the unique replication cost.
%\end{block}
%\end{frame}
%
%\begin{frame}{Can dominance recover the elementary call bounds?}
%For a non-dividend-paying stock,
%\[
%0\leq(S_T-K)^+\leq S_T.
%\]
%\begin{pausebox}
%Translate both payoff inequalities into price inequalities. Then use
%put--call parity and \(P_t\geq0\) for a lower bound.
%\end{pausebox}
%\pause
%The upper bound is immediate:
%\[
%0\leq C_t\leq S_t.
%\]
%\end{frame}
%
%\begin{frame}{Static arbitrage pins the call between two portfolios}
%Put--call parity gives
%\[
%C_t-P_t=S_t-Ke^{-r(T-t)}.
%\]
%Since \(P_t\geq0\),
%\[
%C_t\geq S_t-Ke^{-r(T-t)}.
%\]
%Combining with \(C_t\geq0\),
%\[
%\boxed{
%\bigl(S_t-Ke^{-r(T-t)}\bigr)^+
%\leq C_t\leq S_t.}
%\]
%\begin{marketbox}
%These bounds require no distributional assumption and no option formula.
%\end{marketbox}
%\end{frame}

% ============================================================
\section{Synthesis: Choosing the Right Route}
% ============================================================

\begin{frame}{The PDE and martingale routes answer the same pricing problem}
\begin{columns}[T,onlytextwidth]
\column{0.48\textwidth}
\begin{block}{PDE route}
\begin{itemize}
  \item exploit a Markov state;
  \item apply It\^o and replication;
  \item read delta as a derivative;
  \item solve a differential equation.
\end{itemize}
\end{block}
\column{0.48\textwidth}
\begin{exampleblock}{Martingale route}
\begin{itemize}
  \item choose a numeraire and EMM;
  \item price by conditional expectation;
  \item recover the hedge through MRT;
  \item handle path dependence naturally.
\end{itemize}
\end{exampleblock}
\end{columns}
\end{frame}

\begin{frame}{The payoff structure suggests the most efficient tool}
\begin{center}
\begin{tabular}{p{0.34\textwidth}p{0.50\textwidth}}
\toprule
Problem & Natural first route\\
\midrule
Smooth European payoff in one state & PDE or closed-form expectation\\
Linear payoff such as a forward & replication / law of one price\\
Path-dependent Asian payoff & conditional expectation / Monte Carlo\\
American exercise feature & optimal stopping / Snell envelope\\
Unspanned factor & EMM range / incomplete-market analysis\\
\bottomrule
\end{tabular}
\end{center}
\end{frame}


\begin{frame}{The diagnostic map separates five different ideas}
\begin{enumerate}
  \item Identical payoff \(\Rightarrow\) law of one price.
  \item Untraded Brownian risk \(\Rightarrow\) incompleteness.
  \item Early exercise \(\Rightarrow\) optimal stopping and Snell envelope.
  \item Running average \(\Rightarrow\) state augmentation.
  \item Martingale to holdings \(\Rightarrow\) MRT and exposure matching.
\end{enumerate}
\begin{block}{Common thread}
No-arbitrage restricts prices; spanning determines whether those restrictions
select one price and one attainable payoff.
\end{block}
\end{frame}

\begin{frame}{Use eight checks for every martingale-pricing problem}
\begin{enumerate}
  \item What is the numeraire and discount factor?
  \item Under which measure are discounted traded gains martingales?
  \item Is the payoff replicable, or is some risk unspanned?
  \item Is the relevant local martingale a true martingale?
  \item What state variables summarize the payoff?
  \item Does exercise introduce an optimization over stopping times?
  \item Which Brownian exposures appear in the MRT coefficient?
  \item Do traded assets span those exposures?
\end{enumerate}
\begin{block}{Central message}
Martingale pricing computes values, martingale representation reveals hedges,
and completeness decides whether no-arbitrage selects a unique answer.
\end{block}
\end{frame}

% ============================================================
\section{Revision Exercises: Lectures 2--6}
% ============================================================

\begin{frame}{Revision 1: apply It\^o's formula to a sine}
Let \(X_0=0\) and
\[
\dd X_t=(1+X_t)\dd t+(1+t)\dd W_t,
\qquad Y_t=\sin(X_t).
\]
\begin{pausebox}
Find \(\dd Y_t\), identifying its drift and Brownian coefficient.
\end{pausebox}
\begin{enumerate}
  \item Compute \(f'\) and \(f''\) for \(f(x)=\sin x\).
  \item Compute \(\dd[X]_t\) from the Brownian coefficient of \(X\).
  \item Apply It\^o's formula. Which term would the ordinary chain rule miss?
\end{enumerate}
\end{frame}

\begin{frame}{Revision 1: correction, from derivatives to the SDE}
For \(f(x)=\sin x\), \(f'(x)=\cos x\) and \(f''(x)=-\sin x\).
The Brownian coefficient gives \(\dd[X]_t=(1+t)^2\dd t\).
\[
\begin{aligned}
\dd Y_t
&=f'(X_t)\dd X_t+\tfrac12 f''(X_t)\dd[X]_t\\
&=\cos X_t\bigl[(1+X_t)\dd t+(1+t)\dd W_t\bigr]
 -\tfrac12\sin X_t(1+t)^2\dd t.
\end{aligned}
\]
\[
\boxed{\begin{aligned}
\dd Y_t={}&\left[(1+X_t)\cos X_t
-\tfrac12(1+t)^2\sin X_t\right]\dd t\\
&+(1+t)\cos X_t\dd W_t.
\end{aligned}}
\]
The extra drift \(-\tfrac12(1+t)^2\sin X_t\) comes from quadratic variation.
\end{frame}

\begin{frame}{Revision 2: remove a deterministic drift with Girsanov}
Under \(\Pp\), let \(Y_0=0\) and
\[\dd Y_t=(\alpha+\beta t)\dd t+\dd W_t.\]
\begin{pausebox}
Construct a measure \(\Q\sim\Pp\) on \(\cF_T\) under which \(Y\) is Brownian.
\end{pausebox}
\begin{enumerate}
  \item Choose the drift shift and write its stochastic exponential.
  \item Verify Novikov's condition on the finite horizon \([0,T]\).
  \item Identify \(W^\Q\) and check that it equals \(Y\).
\end{enumerate}
\end{frame}

\begin{frame}{Revision 2: correction, construct a valid density}
Set \(\theta_t=\alpha+\beta t\) and define
\[
Z_t=\exp\!\left(-\int_0^t\theta_s\dd W_s
-\tfrac12\int_0^t\theta_s^2\dd s\right).
\]
Since \(\theta\) is deterministic,
\[
I_T:=\int_0^T\theta_s^2\dd s
=\alpha^2T+\alpha\beta T^2+\frac{\beta^2T^3}{3}<\infty.
\]
\[
\E^\Pp\!\left[e^{\frac12\int_0^T\theta_s^2\dd s}\right]
=e^{I_T/2}<\infty.
\]
Novikov gives \(\E^\Pp[Z_T]=1\). As \(Z_T>0\), it defines
\[
\boxed{\frac{\dd\Q}{\dd\Pp}\bigg|_{\cF_T}=Z_T,\qquad\Q\sim\Pp.}
\]
\end{frame}

\begin{frame}{Revision 2: correction, check the sign of the shift}
Girsanov's theorem says that
\[
\boxed{W_t^\Q=W_t+\int_0^t\theta_s\dd s
=W_t+\alpha t+\tfrac12\beta t^2}
\]
is Brownian under \(\Q\). Hence \(\dd W_t=\dd W_t^\Q-\theta_t\dd t\), and
\[
\dd Y_t=\theta_t\dd t+\dd W_t
=\theta_t\dd t+\dd W_t^\Q-\theta_t\dd t
=\dd W_t^\Q.
\]
Since \(Y_0=W_0^\Q=0\), we have \(Y_t=W_t^\Q\).
\begin{block}{Distribution check}
Under \(\Pp\), \(Y_t\sim\Nlaw(\alpha t+\tfrac12\beta t^2,t)\).
Under \(\Q\), \(Y_t\sim\Nlaw(0,t)\).
\end{block}
\end{frame}

\begin{frame}{Revision 3: the maximum of Brownian motion with drift}
Under \(\Pp\), let \(W\) be Brownian and define
\[
X_t=W_t+\mu t,\qquad M_T^\mu=\max_{0\le s\le T}X_s,
\qquad T>0,\quad a>0.
\]
\begin{pausebox}
Compute \(\Pp(M_T^\mu\ge a)\) using Girsanov and the reflection principle.
\end{pausebox}
\begin{enumerate}
  \item Make \(X\) Brownian under a new measure \(\Q\).
  \item Reflect paths at their first visit to \(a\).
  \item Return to \(\Pp\) and evaluate the Gaussian integrals.
\end{enumerate}
Here \(\Phi\) denotes the standard normal distribution function.
\end{frame}

\begin{frame}{Revision 3: correction, remove the drift first}
Define
\[
\frac{\dd\Q}{\dd\Pp}\bigg|_{\cF_T}
=\exp\!\left(-\mu W_T-\tfrac12\mu^2T\right).
\]
The constant shift satisfies Novikov. Under \(\Q\),
\(X_t=W_t+\mu t\) is a standard Brownian motion.
\[
\frac{\dd\Pp}{\dd\Q}\bigg|_{\cF_T}
=\exp\!\left(\mu X_T-\tfrac12\mu^2T\right).
\]
Therefore the desired probability is
\[
\boxed{\Pp(M_T^\mu\ge a)
=\E^\Q\!\left[e^{\mu X_T-\mu^2T/2}
\1_{\{M_T^\mu\ge a\}}\right].}
\]
The maximum is evaluated on the same paths under both measures.
\end{frame}

\begin{frame}{Revision 3: correction, reflect the paths ending below \(a\)}
Under \(\Q\), \(X_T\) has density
\(\varphi_T(x)=(2\pi T)^{-1/2}e^{-x^2/(2T)}\).
\begin{itemize}
  \item If \(X_T\ge a\), the continuous path has necessarily visited \(a\).
  \item If \(X_T=x<a\) after visiting \(a\), reflection sends its endpoint to
        \(2a-x>a\), preserving Brownian probability.
\end{itemize}
\[
\Q(M_T^\mu\ge a,\ X_T\in\dd x)
=\begin{cases}
\varphi_T(x)\dd x,&x\ge a,\\
\varphi_T(2a-x)\dd x,&x<a.
\end{cases}
\]
Weighting these two parts by \(e^{\mu x-\mu^2T/2}\) gives
\[
\Pp(M_T^\mu\ge a)
=\underbrace{\int_a^\infty e^{\mu x-\mu^2T/2}\varphi_T(x)\dd x}_{I_1}
+\underbrace{\int_{-\infty}^a e^{\mu x-\mu^2T/2}
\varphi_T(2a-x)\dd x}_{I_2}.
\]
\end{frame}

\begin{frame}{Revision 3: correction, evaluate the two Gaussian tails}
Completing the square gives
\(e^{\mu x-\mu^2T/2}\varphi_T(x)=\varphi_T(x-\mu T)\), so
\[
I_1=\int_a^\infty\varphi_T(x-\mu T)\dd x
=\Phi\!\left(\frac{\mu T-a}{\sqrt T}\right).
\]
For \(I_2\), substitute \(y=2a-x\):
\[
I_2=e^{2\mu a}\int_a^\infty e^{-\mu y-\mu^2T/2}\varphi_T(y)\dd y
=e^{2\mu a}\Phi\!\left(\frac{-\mu T-a}{\sqrt T}\right).
\]
\[
\boxed{\Pp(M_T^\mu\ge a)
=\Phi\!\left(\frac{\mu T-a}{\sqrt T}\right)
+e^{2\mu a}\Phi\!\left(\frac{-\mu T-a}{\sqrt T}\right).}
\]
For \(\mu=0\), we recover \(2[1-\Phi(a/\sqrt T)]\).
\end{frame}

\begin{frame}{Revision 4: which random times are stopping times?}
Use the Brownian filtration \((\cF_t)\). Fix \(T>0\) and \(a>0\).
\begin{pausebox}
For each random time, decide whether \(\{\tau\le t\}\in\cF_t\) for every \(t\).
\end{pausebox}
\begin{enumerate}
  \item \(\tau_a=\inf\{t\ge0:W_t\ge a\}\);
  \item \(T/2\);
  \item \(\eta=\inf\{t\ge T/2:W_t=0\}\);
  \item \(\rho_T=\sup\{0\le t\le T:W_t=0\}\);
  \item \(\gamma_T=\operatorname*{arg\,max}_{0\le s\le T}W_s\).
\end{enumerate}
The Brownian maximum has a unique maximizing time almost surely.
\end{frame}

\begin{frame}{Revision 4: correction, the first three are stopping times}
\textbf{First visit to \(a\).} By continuity,
\[
\{\tau_a\le t\}=\left\{\max_{0\le s\le t}W_s\ge a\right\}\in\cF_t.
\]
\textbf{The deterministic time \(T/2\).} Its event is
\[
\{T/2\le t\}=\begin{cases}\varnothing,&t<T/2,\\\Omega,&t\ge T/2.\end{cases}
\]
\textbf{First zero after \(T/2\).} Again by continuity,
\[
\{\eta\le t\}=\begin{cases}
\varnothing,&t<T/2,\\
\left\{\min_{T/2\le s\le t}|W_s|=0\right\},&t\ge T/2.
\end{cases}
\]
Each event can be decided from the path observed up to \(t\).
\end{frame}

\begin{frame}{Revision 4: correction, the last zero needs future information}
Fix \(0<t<T\). The event is
\[
\{\rho_T\le t\}=\{W\text{ has no zero in }(t,T]\}.
\]
Knowing \(W_t\neq0\) does not tell us whether a future return to zero occurs.
By independent increments and the reflection principle,
\[
\boxed{\Pp(\rho_T\le t\mid\cF_t)
=2\Phi\!\left(\frac{|W_t|}{\sqrt{T-t}}\right)-1.}
\]
For a fixed \(0<t<T\), this lies strictly between \(0\) and \(1\) almost surely.
\begin{block}{Why this proves nonmeasurability}
If \(\{\rho_T\le t\}\) belonged to \(\cF_t\), its conditional probability
would be its indicator, hence either \(0\) or \(1\).
Thus \(\rho_T\) is not a stopping time.
\end{block}
\end{frame}

\begin{frame}{Revision 4: correction, the maximizing time also looks ahead}
Fix \(0<t<T\) and let \(m_t=\max_{0\le s\le t}W_s\).
Up to a null event,
\[
\{\gamma_T\le t\}
=\left\{\max_{t\le s\le T}W_s\le m_t\right\}.
\]
The past tells us \(m_t\) and \(W_t\), but a future increment may exceed
the remaining gap \(m_t-W_t\). Reflection gives
\[
\boxed{\Pp(\gamma_T\le t\mid\cF_t)
=2\Phi\!\left(\frac{m_t-W_t}{\sqrt{T-t}}\right)-1.}
\]
\begin{block}{The same measurability test}
For a fixed \(0<t<T\), the gap is positive almost surely, and the displayed
probability lies in \((0,1)\). It cannot be the indicator of an
\(\cF_t\)-measurable event. Thus \(\gamma_T\) is not a stopping time.
\end{block}
\end{frame}

\begin{frame}{Final recap: quadratic variation and covariation}
For partitions \(\pi=(t_k)\) of \([0,t]\), put
\(\Delta X_k=X_{t_{k+1}}-X_{t_k}\). As the mesh tends to zero,
\[
\sum_k(\Delta X_k)^2\xrightarrow{\Pp}[X]_t,
\qquad
\sum_k\Delta X_k\Delta Y_k\xrightarrow{\Pp}[X,Y]_t.
\]
\begin{block}{What survives for continuous It\^o processes}
If \(\dd X_t=a_t\dd t+\sum_i b_t^i\dd W_t^i\) and
\(\dd Y_t=c_t\dd t+\sum_i d_t^i\dd W_t^i\), with independent Brownian
motions, then
\[
\dd[X]_t=\sum_i(b_t^i)^2\dd t,
\qquad
\dd[X,Y]_t=\sum_i b_t^i d_t^i\dd t.
\]
\end{block}
In particular, \([W]_t=t\); finite-variation terms contribute zero.
\end{frame}

\begin{frame}{Final recap: the algebra of stochastic differentials}
\[
(\dd t)^2=0,\qquad \dd t\,\dd W_t^i=0,\qquad
\dd W_t^i\,\dd W_t^j=\delta_{ij}\dd t
\quad\text{(independent Brownian motions).}
\]
Read \(\dd X_t\,\dd Y_t\) as \(\dd[X,Y]_t\). Thus
\[
\boxed{\dd(X_tY_t)=X_t\dd Y_t+Y_t\dd X_t+\dd X_t\,\dd Y_t.}
\]
\begin{pausebox}
If \(\dd X_t=2\dd t+2\dd W_t^1+\dd W_t^2\) and
\(\dd Y_t=-\dd t-\dd W_t^1+3\dd W_t^2\), what is
\(\dd X_t\,\dd Y_t\)?
\end{pausebox}
\end{frame}

\begin{frame}{Final correction: expand the product of the increments}
All products containing \(\dd t\) vanish, leaving
\[
\begin{aligned}
\dd X_t\,\dd Y_t
&=(2\dd W_t^1+\dd W_t^2)(-\dd W_t^1+3\dd W_t^2)\\
&=-2(\dd W_t^1)^2+6\dd W_t^1\dd W_t^2
 -\dd W_t^2\dd W_t^1+3(\dd W_t^2)^2\\
&=(-2+3)\dd t=\boxed{\dd t}.
\end{aligned}
\]
The mixed Brownian products vanish because \(W^1\) and \(W^2\) are independent.
Applying the product rule now gives
\[
\boxed{\begin{aligned}
\dd(X_tY_t)={}&(2Y_t-X_t+1)\dd t\\
&+(2Y_t-X_t)\dd W_t^1+(Y_t+3X_t)\dd W_t^2.
\end{aligned}}
\]
\end{frame}


\end{document}
